\documentclass[10pt,a4paper]{amsart}
\usepackage[margin=19mm]{geometry}
\usepackage{amsmath,amssymb,mathtools}
\usepackage[scr=rsfs,cal=euler]{mathalfa}
\usepackage{xfrac}
\usepackage[unicode,hidelinks,bookmarks=false]{hyperref}
\newcommand{\ie}{i.e.\ }
\newcommand{\cf}{cf.\ }
\newcommand{\resp}{respectively\ }
\newcommand{\Subsection}{\subsection}
\providecommand{\nobreakdash}{\nobreak\hbox{-}}
\setlength{\emergencystretch}{2em}
\multlinegap0pt
\usepackage{bm}
\usepackage{bbm}
\usepackage{dsfont}
\usepackage{xspace}
\usepackage{cancel}
\usepackage{mathtools}
\usepackage{amsthm}
\makeatletter
\@ifundefined{theorem}{\newtheorem{theorem}{Theorem}}{}
\@ifundefined{lemma}{\newtheorem{lemma}[theorem]{Lemma}}{}
\makeatother
\title[Circle packings and hyperbolic surfaces of finite type]{Circle packings and hyperbolic surfaces of finite type}
\author{Te Ba, Guangming Hu, Yu Sun}
\date{Source: arXiv:2307.13572v3 (CC BY 4.0)}
\begin{document}
\maketitle

\noindent\emph{Source and licence.} Circle packings and hyperbolic surfaces of finite type. Version arXiv:2307.13572v3 (CC BY 4.0), distributed under the Creative Commons Attribution 4.0 International licence, as recorded in the arXiv licence metadata for this exact version and shown on the abstract page https://arxiv.org/abs/2307.13572v3. Full rights notice: licenses/arxiv-2307.13572-CC-BY-4.0.txt.\par\smallskip

\noindent\textit{Editorial scope.} Counted unit: the Lemma whose source label is \texttt{L-3-5} in the article (source0.tex lines 458--466, source label \texttt{L-3-5}), delivered together with the article's own complete original proof of it (source0.tex lines 467--505, one \texttt{proof} environment of 39 lines). No mathematics is written, repaired or re-derived: statement and proof are the article's own text line for line. Required context retained uncounted: E-gh (lines 301--301); E-2-2 (lines 303--314); E-2-3 (lines 316--319); L-3-3 (lines 329--336); E-3-7 (lines 379--379); E-2-8 (lines 384--384); E-3-1 (lines 410--410); E-2-9 (lines 433--435); E-2-10 (lines 437--440); E-6-1 (lines 733--738); E-6-3 (lines 761--761). Every definition, display and label of those lines is retained unchanged. Omitted from this package and not redistributed: 1--300, 302--302, 315--315, 320--328, 337--378, 380--383, 385--409, 411--432, 436--436, 441--457, 506--732, 739--760, 762--975. Nothing inside a retained excerpt is omitted or altered: every retained body is delivered exactly as the article's lines are written, so no omission receipt of any kind accompanies this package. Citations: the retained excerpts carry no citation command at all, so no bibliography entry of the article is transcribed and no reference is redistributed. Reference adaptation: none was needed and none was made; every reference inside the delivered unit resolves inside it, so no unresolved reference or citation is produced. Printed slips kept literal: no printed word, symbol, hypothesis or number of the counted statement or of its proof was altered. Layout and class: the article's own class is \texttt{article} with packages including babel, tikz, graphicx, booktabs, tabularx, verbatim, etoolbox, tkz-euclide, geometry, mathtools, color, appendix, xy, tikz-cd; the wrapper is the lane's pinned \texttt{amsart} editorial wrapper at 10pt on A4, which restates the theorem-like environments the retained text needs and adds the article's own macro definitions that the retained text needs (none), with extra wrapper packages (bm, bbm, dsfont, xspace, cancel, mathtools). The delivered numbering is the unit's own. Assets: no retained span contains a figure environment, a graphics-inclusion command, a PDF-page inclusion, a table or a TikZ picture; no image file is copied into this package, the package declares no asset dependency and no image file is redistributed with it. Licence: version arXiv:2307.13572v3 of this article is distributed under the Creative Commons Attribution 4.0 International licence, as recorded in the arXiv licence metadata for this exact version and shown on the abstract page https://arxiv.org/abs/2307.13572v3. A full rights notice is delivered at licenses/arxiv-2307.13572-CC-BY-4.0.txt. Characters: every retained excerpt is pure ASCII, so the delivered engine, XeLaTeX with its default Latin Modern text font, reports no missing character.
\begin{equation}\label{E-gh}g(\kappa_1,\kappa_2)=\int_{\mathcal C_1}\kappa_1ds,\quad h(\kappa_1,\kappa_2)=\int_{\mathcal C_2}\kappa_2ds.\end{equation}

\begin{equation}\label{E-2-2}
g(\kappa_1,\kappa_2)
=
\begin{cases}
2b_{\kappa_1}\kappa_1\cot^{-1}(b_{\kappa_1}\kappa_2),
& \kappa_1>1,\\
\frac{2\kappa_1}{\kappa_2},
& \kappa_1=1,\\
2b_{\kappa_1}\kappa_1\coth^{-1}(b_{\kappa_1}\kappa_2),
& \kappa_1<1.
\end{cases}
\end{equation}

\begin{equation}\label{E-2-3}
h(\kappa_1,\kappa_2)
=2b_{\kappa_2}\kappa_2\cot^{-1}(b_{\kappa_2}\kappa_1).
\end{equation}

\begin{lemma}\label{L-3-3}
 Let $\mathcal C_1$ and $\mathcal C_2$ be two sides of a right-angled bigon in~$\mathbb H^2$ with curvatures $\kappa_1,$ and $\kappa_2$ respectively, and assume that $\kappa_2\geq 1$. Let $g$, $h$ be the total geodesic curvatures of $\mathcal C_1$ and $\mathcal C_2$ as defined in (\ref{E-gh}).
 Then the following statements hold:
\begin{itemize}
\item[($i$)] $g$, $h$ are given by (\ref{E-2-2}) and (\ref{E-2-3}).
\item[($ii$)]  $g$, $h$ are $C^1$-smooth functions.
\end{itemize}
\end{lemma}

\begin{equation}\label{E-3-7}\theta(\kappa_s,\kappa)=\frac{h(\kappa_s,\kappa)}{\cosh r}=2\cot^{-1}(b_{\kappa_s}\kappa)\end{equation}

\begin{equation}\label{E-2-8}\frac{d\theta(\kappa_s,k)}{d\kappa}=\frac{2\kappa_s\kappa b_{\kappa}}{\kappa_s^2+\kappa^2-1}>0,\end{equation}

\begin{equation}\label{E-3-1}f_i(\kappa_u,\kappa_v,\kappa_w)=\int_{\mathcal C_i}\kappa_ids\end{equation}

\begin{equation}\label{E-2-9}\
g_y(\kappa_1,\kappa_2)=\frac{2\kappa_1}{1-\kappa_1^2-\kappa_2^2}.
\end{equation}

\begin{equation}\label{E-2-10}\begin{aligned}
\frac{\partial}{\partial k_v}f_u(\kappa_u,\kappa_v,\kappa_w)&=g_y(\kappa_u,\kappa_c)\frac{\partial \kappa_c}{\partial \kappa_v}\frac{d\kappa_v}{dk_v}&=\frac{2\kappa_u\kappa_v}{1-\kappa_u^2-\kappa_c^2}\frac{\kappa_u+\kappa_w}{2\kappa_c},\\
\frac{\partial}{\partial k_u}f_v(\kappa_u,\kappa_v,\kappa_w)&=g_y(\kappa_v,\kappa_c)\frac{\partial \kappa_c}{\partial \kappa_u}\frac{d\kappa_u}{dk_u}&=\frac{2\kappa_u\kappa_v}{1-\kappa_v^2-\kappa_c^2}\frac{\kappa_v+\kappa_w}{2\kappa_c}.
\end{aligned}\end{equation}

\begin{equation}\label{E-6-1}\begin{aligned}
g_x(\kappa_1,\kappa_2)=\begin{cases}
\frac{2\kappa_1^2\kappa_2b^2_{\kappa_1}}{\kappa_1^2+\kappa_2^2-1}-2b_{\kappa_1}^3\cot^{-1}b_{\kappa_1}\kappa_2, & \kappa_1>1,\\[1em]
\frac{-2\kappa_1^2\kappa_2b^2_{\kappa_1}}{\kappa_1^2+\kappa_2^2-1}+2b_{\kappa_1}^3\coth^{-1}b_{\kappa_1}\kappa_2, & 0<\kappa_1<1,\\
\end{cases}
\end{aligned}\end{equation}

\begin{equation}\label{E-6-3}\lim_{\kappa_1\to 1}g_x(\kappa_1,\kappa_2)=\frac{2}{\kappa_2}-\frac{4}{3\kappa^3_2}.\end{equation}

\begin{lemma}\label{L-3-5}
Let $\mathcal C_u,\mathcal C_v,\mathcal C_w: [0,1]\to\mathbb{H}^2$ be three smooth arcs of constant curvatures $\kappa_u, \kappa_v, \kappa_w>0$, respectively, such that each pair is externally tangent at their common endpoint.
Let $f_i$ be defined as in (\ref{E-3-1}). Then
\begin{subequations}
\begin{align}
\label{E-3-9a}\frac{\partial f_u}{\partial \kappa_v}&<0,\\
\label{E-3-9b}\frac{\partial(f_u+f_v+f_w)}{\partial \kappa_u}&>0.
 \end{align}\end{subequations}
\end{lemma}

\begin{proof}
We readily obtain (\ref{E-3-9a}) from (\ref{E-2-10}). Let us prove (\ref{E-3-9b}).
We use $\kappa_c$ to denote the curvature passing through three points of tangency.
To simplify the computations, we derive the substitution
\begin{equation}\label{E-3-10}
\frac{\partial \kappa_c}{\partial \kappa_u}=\frac{-\theta_x(k_u,k_c)}{\theta_y(k_u,k_c)+\theta_y(k_v,k_c)+\theta_y(k_w,k_c)},
\end{equation}
which arises from differentiating
 \[\theta(k_u,k_c)+\theta(k_v,k_c)+\theta(k_w,k_c)=2\pi\]
on both sides with respect to $\kappa_u$. Meanwhile, we compute from (\ref{E-3-7}) that
\begin{equation}\label{E-3-3}\theta_x(\kappa_u,\kappa_c)=\frac{2}{b_{\kappa_c}(1-\kappa_u^2-\kappa_c^2)}.\end{equation}
Comparing (\ref{E-2-8}) and (\ref{E-2-9}), we observe that
\begin{equation}\label{E-3-4}\theta_y(\kappa_i,\kappa_c)=-g_y(\kappa_i,\kappa_c)\kappa_cb_{\kappa_c}\end{equation}
for $i\in\{u,v,w\}$. According to (\ref{E-3-10}), (\ref{E-3-3}) and (\ref{E-3-4}), we deduce that
 \[\begin{aligned}
 \frac{\partial}{\partial \kappa_u}(f_u+f_v+f_w)&=\frac{\partial}{\partial k_u}(g(\kappa_u,\kappa_c)+g(\kappa_v,\kappa_c)+g(\kappa_w,\kappa_c))\frac{\partial \kappa_c}{\partial \kappa_u}\\
 &=g_x(\kappa_u,\kappa_c)+\big(g_y(\kappa_u,\kappa_c)+g_y(\kappa_v,\kappa_c)+g_y(\kappa_w,\kappa_c)\big)\frac{\partial \kappa_c}{\partial \kappa_u}\\
 &=g_x(\kappa_u,\kappa_c)+\frac{\theta_x(\kappa_u,\kappa_c)}{\kappa_cb_{\kappa_c}}\\
&=g_x(\kappa_u,\kappa_c)+\frac{2}{\kappa_c(1 - \kappa_u^2 - \kappa_c^2)b^2_{\kappa_c}}.
 \end{aligned}\]
We now apply the result of  Lemma \ref{L-3-3}($ii$), (\ref{E-6-1}) and (\ref{E-6-3}) to obtain



  \[
\frac{\partial}{\partial k_u}(f_u+f_v+f_w)
=
\begin{cases}
\frac{2}{(k_u^2-1)k_c}-2b^3_{k_u}\cot^{-1}b_{k_u}k_c,
& k_u>1,\\
\frac{2}{3k_c^3}
& k_u=1,\\
\frac{2}{(k_u^2-1)k_c}+2b^3_{k_u}\cot^{-1}b_{k_u}k_c,
& k_u<1.
\end{cases}
\]
Observe that $\cot^{-1}b_{k_u}k_c<\frac{1}{b_{k_u}k_c}$ when $k_u\neq 1$. It follows that
\[\frac{\partial}{\partial k_u}(f_u+f_v+f_w)>0.\]
\end{proof}

\end{document}
